% Figure: position, velocity, acceleration versus time for the two-phase
% car motion (accelerate then decelerate). Stacked panels share a time axis.
\begin{figure}[htbp]
\centering
\begin{tikzpicture}
\begin{axis}[
  name=accax,
  width=11cm, height=3.2cm,
  ylabel={$a$ (\si{\meter\per\second\squared})},
  xmin=0, xmax=25, ymin=-2.5, ymax=3.5,
  axis lines=left, xtick=\empty,
  grid=major, grid style={enggray!20},
  tick label style={font=\scriptsize}, label style={font=\small},
]
% phase 1: a = +3 for 0..10; phase 2: a = -2 for 10..25
\addplot[engred, very thick] coordinates {(0,3) (10,3)};
\addplot[engred, very thick] coordinates {(10,-2) (25,-2)};
\addplot[engred, very thick, dashed] coordinates {(10,3) (10,-2)};
\end{axis}
\begin{axis}[
  name=velax, at={(accax.below south west)}, anchor=above north west, yshift=-0.3cm,
  width=11cm, height=3.2cm,
  ylabel={$v$ (\si{\meter\per\second})},
  xmin=0, xmax=25, ymin=0, ymax=35,
  axis lines=left, xtick=\empty,
  grid=major, grid style={enggray!20},
  tick label style={font=\scriptsize}, label style={font=\small},
]
% v rises 0->30 over 0..10, falls 30->0 over 10..25
\addplot[engblue, very thick] coordinates {(0,0) (10,30) (25,0)};
\end{axis}
\begin{axis}[
  name=posax, at={(velax.below south west)}, anchor=above north west, yshift=-0.3cm,
  width=11cm, height=3.2cm,
  xlabel={time $t$ (\si{\second})}, ylabel={$x$ (\si{\meter})},
  xmin=0, xmax=25, ymin=0, ymax=400,
  axis lines=left,
  grid=major, grid style={enggray!20},
  tick label style={font=\scriptsize}, label style={font=\small},
]
% x: integral of v. Phase 1: x=1.5 t^2 (0..10 -> 150). Phase 2: x=150+30(t-10)-1(t-10)^2.
\addplot[engamber, very thick, domain=0:10, samples=40] {1.5*x^2};
\addplot[engamber, very thick, domain=10:25, samples=60] {150 + 30*(x-10) - 1*(x-10)^2};
\end{axis}
\end{tikzpicture}
\caption[Stacked kinematics of two-phase motion]{Acceleration,
velocity, and position against a shared time axis for the
two-phase car of the second calculation exercise: constant
$+3\,\si{\meter\per\second\squared}$ for the first
$10\,\si{\second}$, then constant $-2\,\si{\meter\per\second\squared}$
to rest at $t=25\,\si{\second}$. Velocity is the running integral of
acceleration, so the step in $a$ becomes a kink in $v$; position is
the running integral of velocity, so the kink in $v$ becomes a change
of curvature in $x$. The area under the velocity triangle,
$\tfrac{1}{2}(25)(30)=375\,\si{\meter}$, is the total distance, a
reading that the integral confirms. Reading these three panels
together is the core skill of one-dimensional kinematics.}
\label{fig:vol03ch04:kinematics-1d}
\end{figure}
